\section{Introduction}
\label{sec:intro}
视觉模仿学习通过示范建立观察到动作的映射，使机器人能够学习难以显式编程的操作。
动作块策略进一步预测一段未来动作，为连续操作提供计划上下文。
然而，在机器人示范数量有限时，拟合示范动作并不意味着可靠完成任务：
工具接触前的局部偏差可能导致拨空，而大量自由运动帧上的准确预测未必能补偿这一失败。
本文关注的问题是：\emph{在有限示范条件下，如何通过语义引导的任务关系表示与历史条件反馈，
提高动作块模仿策略的任务完成表现？}

这一问题包含两个互补缺口。首先，有限动作监督未必足以使视觉表示保留关键控制关系。
例如，在遮挡物移开后推动目标进入区域，策略需要区分目标是否暴露、工具相对目标的位置以及目标与区域的关系。
这些关系虽然存在于RGB中，却不一定能仅凭少量动作示范被充分学到。
语义监督提供对象类别与空间范围，关系监督进一步要求表示保留与操作有关的几何量。
它们提供的是有限数据下的学习先验，而不是给已有RGB增加新的测量信息；
分割错误、关系选择不当或动作预测未利用这些关系，都可能抵消收益。

其次，准确表示生成动作块时的画面，仍不能保证执行期间的信息充分。
当工具接近目标时，目标可能移动或被遮挡；相似的当前几何关系也可能对应接近、接触或滑脱等不同动态状态。
新观测用于发现原计划形成后的变化，历史用于区分单帧难以确定的动态上下文。
动作块策略可以通过重规划或时间集成使用新观测，并非完全开环；
但这些机制不自动等同于利用观察历史、针对当前计划学习局部修正。
本文因此研究计划条件反馈，而不预设提高重规划频率或使用循环网络就必然改善表现。
\draftnote{待补动机图：用可追溯的执行序列标出观察、基础计划、接触变化与修正目标；未验证的失效原因仅作为待检验解释。}

针对表示与执行信息两个缺口，本文构建语义引导的任务关系表示与历史条件残差反馈框架。
基础策略结合RGB、语义图及几何监督查询生成动作块；冻结基础策略后，
反馈分支根据新观测、因果历史和当前计划修正尚未执行的动作槽位。
两条路径服务于同一任务完成目标，但作用不同：表示监督不能补出未来变化，
历史反馈也不能自动保证识别正确的对象关系。
本文的核心待验证主张是，二者在表示可靠、修正有效且时序对齐的条件下，
能够共同改善有限示范训练的动作块策略表现，而不是对任意任务给出无条件保证。

本文围绕三个方面组织工作：
\begin{enumerate}
\item 建立从任务相关动作误差、表示误差到历史反馈收益的条件性分析，
明确各环节的适用条件与可检验量；已知性能差分和投影恒等式作为分析工具，而非新的基础定理。
\item 构建语义与关系监督的动作块策略及计划条件历史反馈，
通过目标槽位对齐将局部修正作用于对应的未来动作。
\item 设计表示与反馈的交叉对照、示范规模实验和机制诊断，
分别检验任务表现、两条路径的贡献以及反馈失效条件。
\end{enumerate}
\draftnote{待补：以上实验设计尚不构成已完成的实验证据；根据最终核验结果改写实验贡献并填入核心定量结论。}
本文中的有限示范指机器人操作示范受限，不等于总监督资源受限；
额外分割标签、预训练及人工审核成本均需披露。
单一示范规模的改善只能支持该设置下的有效性，样本效率主张需要多个规模的比较。
第二节定位相关工作，第三节给出问题定义与条件性分析，第四节描述方法，
第五节围绕研究问题组织验证，第六节讨论适用边界与结论。
